\section{Introduction}
\label{sec:introduction}

An approximate attention value and a certified output decision solve different
problems. A small absolute error can still cross a rounding boundary, while
a coarse value estimate may determine a broad downstream decision exactly.
The relevant quantity is therefore the uncertainty relative to the consumer's
decision boundary.

We study positive normalized reductions represented by block masses and
centered value moments. At any boundary $c$, the normalized comparison reduces
to the sign of an unnormalized residual $N-cZ$. Independent interval metadata
gives attained extrema of this residual. If total mass is positive throughout
the box, two strict sign tests are necessary and sufficient for every compatible
output to lie in a given open cell. A rejection identifies information missing from the box;
it need not identify an inaccurate candidate value.

We test which dependence information enables certification, which error term
prevents it on trained activations, and whether avoided token work repays
construction, screening, and fallback. The evidence combines value-coupled
bounds, a formal real-attention proof chain, and GH200 implementation ablations.

Value coupling adds certificates, although a global value hull explains most
gains. Trained activations expose discarded-score uncertainty and full-rank
Taylor inflation. Parallel reduction removes a GPU bottleneck, yet the
complete path remains slower than fused dense attention throughout the sweep.

\paragraph{Scope.}
Lean proves the real-arithmetic chain under concrete centering, radius,
symmetry, and row-sum witnesses. The rational checker and the directed CUDA
checker are tested implementations, not extracted proofs. Ordinary
NumPy/CUDA summaries do not provide outward intervals. Exact-real rounding
also differs from the accumulation contract of a particular fused GPU kernel.
\Cref{app:formalization} maps these claims to the checked theorems and remaining
implementation bridges. Code, measurements, and reproduction commands are
available in the \href{https://github.com/SamMausberg/contracted-moment-kernels}{research repository}.

\paragraph{Related work.}
Adaptive predicates refine uncertain signs \citep{shewchuk}, while Taylor
models retain polynomial structure with bounded remainders \citep{berz}.
Multipole attention combines grouped interactions and selected exact work
\citep{multipole,fma}; TaylorShift and symmetry-aware features reorganize
polynomial attention \citep{taylorshift,sata}. COBS uses covariance statistics
for block selection, SPLA combines sparse and residual linear paths, and
LOCKS studies limits of moment summaries for broad scores
\citep{cobs,spla,locks}. WitCert and attention-memory observability contracts
already connect runtime witnesses, formal scope, and fallback
\citep{observability,witcert}. Our focus is the consumer-specific residual
interface, the projected signed moments supplying it, and stronger mass/value
dependence. These general principles are established; comparisons with
unreproduced systems concern method scope, not measured performance.
